\documentclass[10pt,letterpaper]{article}
\usepackage[utf8]{inputenc}
\usepackage[T1]{fontenc}
\usepackage[spanish,es-tabla]{babel}
\usepackage[margin=1.8cm]{geometry}
\usepackage{booktabs,array,graphicx,amsmath,hyperref,enumitem,xcolor}
\usepackage{lmodern}
\setlist{nosep,leftmargin=1.2em}
\setlength{\parskip}{2pt}
\renewcommand{\arraystretch}{1.05}
\newcommand{\todo}[1]{\textcolor{red}{\textbf{[PENDIENTE: #1]}}}
\hypersetup{colorlinks=true,urlcolor=blue}

\title{\vspace{-1.2cm}\textbf{CC3092 -- Laboratorio 9: Series de tiempo y forecasting}\
\large Consumo eléctrico por hora (ElectricityLoadDiagrams20112014), horizonte $H=24$ h}
\author{Ángel Esquit \and \todo{integrante 2, carnés}}
\date{\vspace{-0.8cm}Universidad del Valle de Guatemala -- 1 de noviembre de 2026}

\begin{document}
\maketitle

\section{Preparación de la serie}
\textbf{Datos.} Los 140{,}256 registros de 15 min (kW) se convirtieron a kWh ($\times 0.25$) y se sumaron por hora.
Tras filtrar el periodo 2012--2014 y exigir datos completos desde el 2012-01-01 quedan \textbf{320 de 370 clientes},
con \textbf{26{,}304 horas} por serie y \textbf{8{,}417{,}280} observaciones en el panel.
Partición cronológica: entrenamiento hasta 2014-06-30, validación hasta 2014-09-30 y prueba 2014-10-01 a 2014-12-31.

\textbf{Cambio de horario.} En marzo la hora 01:00 no existe (queda con $\approx4\%$ del consumo normal): se marca como
faltante y se interpola linealmente. En octubre la hora 01:00 acumula dos horas reales ($\approx2\times$): se reemplaza por el
promedio de la misma hora $\pm7$ días.

\textbf{Escalas.} El consumo medio por cliente (entrenamiento) va de 2.4 a 48{,}817~kWh (mediana 147; razón máx/mín $\approx 2\times10^4$).
Una normalización global dejaría que unos pocos clientes dominen la pérdida, por lo que se usa \textbf{z-score por serie con
media y desviación calculadas solo en entrenamiento}; las métricas siempre se reportan en kWh.

\textbf{Calendario.} Hora, día de la semana y mes con seno/coseno (preserva la ciclicidad con 2 dimensiones por variable,
en vez de 24+7+12 del one-hot) y una binaria de feriado de Portugal. Todas son deterministas, así que se conocen a futuro.
\todo{describir perfiles hora/día y efecto de Navidad/Año Nuevo con la figura}\;
Consumo medio global: 636~kWh; 24--26 dic: 477; 31 dic--2 ene: 489 (\,$\approx-25\%$).

\textbf{Ventanas.} Con $L=168$, $H=24$, stride 1 h el \texttt{Dataset} produce
\textbf{6{,}943{,}040} ventanas de entrenamiento, \textbf{699{,}200} de validación y \textbf{699{,}200} de prueba (320 series).

\textbf{Auditoría de fugas.} (1) \emph{Split aleatorio}: el corte es cronológico y se verifica con \texttt{assert} que ningún
target de un split cae en el siguiente. (2) \emph{Escalado}: $\mu,\sigma$ solo de entrenamiento. (3) \emph{Early stopping}:
solo con validación; la prueba se evalúa una vez por modelo. (4) \emph{Covariables}: la entrada es $[t-L,t-1]$, el target
$[t,t+H-1]$ y solo el calendario determinista se usa a futuro (verificado con \texttt{assert} sobre una ventana).

\section{Investigación de herramientas}
\begingroup\small\sloppy
\textbf{Herramientas de PyTorch.}
\begin{itemize}
\item \texttt{nn.LSTM / nn.GRU}: \texttt{input\_size} (features por paso), \texttt{hidden\_size}, \texttt{num\_layers},
\texttt{batch\_first} (si es \texttt{True} la entrada es $(B,L,F)$) y \texttt{dropout} (entre capas, no tras la última;
solo actúa si \texttt{num\_layers}$>1$). Salidas: \texttt{output} $(B,L,D\cdot h)$ con el estado oculto de la última capa en
cada paso; \texttt{h\_n} $(D\cdot\text{capas},B,h)$ con el último estado de cada capa; \texttt{c\_n}, de igual forma, solo en
la LSTM (celda de memoria; la GRU no la tiene). $D=2$ si es bidireccional.
\item \texttt{nn.Conv1d} con \texttt{dilation}: se hace causal rellenando $(k-1)d$ ceros \emph{solo a la izquierda}
(\texttt{F.pad}) y usando \texttt{padding=0}, de modo que la salida en $t$ no ve $t+1$. El campo receptivo de una TCN es
$RF=1+\sum_{i}(k-1)\,d_i$ (se duplica el término si cada bloque residual usa dos convoluciones). Con $d_i=2^i$ y $k=3$
(una conv por nivel), $RF=1+2(2^n-1)$: hacen falta $n=7$ niveles ($RF=255$) para cubrir $L=168$; con $n=6$ solo se ven 127 h.
\item \texttt{nn.TransformerEncoderLayer}: \texttt{d\_model} (dimensión de embedding), \texttt{nhead} (cabezas; \texttt{d\_model}
debe ser divisible por \texttt{nhead}), \texttt{dim\_feedforward} (capa oculta de la MLP interna, 2048 por defecto), además de
\texttt{dropout}, \texttt{batch\_first} y \texttt{norm\_first}. La atención no tiene noción de orden, por lo que se \emph{suma}
al embedding de entrada una codificación posicional (seno/coseno fija, $PE_{(p,2i)}=\sin(p/10^{4\cdot 2i/d})$, o aprendida con
\texttt{nn.Embedding}); PatchTST, además, agrupa la serie en parches antes de la atención. Para pronóstico autorregresivo se pasa una
máscara causal.
\end{itemize}
\textbf{Librerías de forecasting.}
\begin{itemize}
\item \textbf{StatsForecast / NeuralForecast} (Nixtla): entrada en formato largo con columnas \texttt{unique\_id}, \texttt{ds}, \texttt{y}
(más columnas exógenas). \texttt{cross\_validation} hace backtesting de origen móvil: \texttt{n\_windows} = número de
orígenes (cutoffs), \texttt{step\_size} = horas entre un origen y el siguiente (24 en este lab) y \texttt{refit}: si es
\texttt{True} reajusta el modelo en cada ventana; si es \texttt{False} entrena una sola vez y reutiliza los pesos (en NeuralForecast
admite también un entero, para reentrenar cada $k$ ventanas). Devuelve \texttt{cutoff}, \texttt{y} y una columna por modelo.
\item \textbf{Chronos}: \texttt{BaseChronosPipeline.from\_pretrained} (modelos \texttt{chronos-bolt}); \texttt{predict\_quantiles(context,
prediction\_length, quantile\_levels)} recibe el contexto (tensor 1D/2D o lista de tensores de distinta longitud), el horizonte y los
cuantiles; devuelve \texttt{(quantiles, mean)} con \texttt{quantiles} de forma $(B,H,|q|)$. Es zero-shot: no se entrena.
\end{itemize}
\textbf{Conceptos.}
\begin{itemize}
\item \emph{Local vs global}: local = un modelo por serie (se adapta a cada una, pero con pocos datos y costo $\propto N$);
global = un modelo para todas (comparte patrones y generaliza mejor, requiere normalizar o identificar la serie).
\item \emph{Teacher forcing}: en entrenamiento el decoder recibe el valor \emph{real} del paso anterior en vez de su propia
predicción, lo que acelera la convergencia pero crea \emph{exposure bias}: en inferencia usa sus predicciones y los errores se acumulan.
\item \emph{RevIN}: normaliza cada ventana con su media y desviación (con parámetros afín aprendibles) antes del modelo y
revierte la normalización a la salida; ayuda ante cambios de nivel (drift) porque la red ve siempre series de escala similar.
\item \emph{Pinball loss} para el cuantil $\tau$: $\ell_\tau(y,\hat q)=\max\{\tau(y-\hat q),(\tau-1)(y-\hat q)\}$; su minimizador es
el cuantil $\tau$ de la distribución, y promediar sobre varios $\tau$ aproxima el CRPS.
\item \emph{MASE vs WAPE}: ambas son libres de escala. MASE divide el MAE entre el MAE del naive estacional en entrenamiento
de \emph{cada serie}, y luego promedia: pondera igual a todos los clientes y es $<1$ si se supera al naive.
WAPE $=\sum|e|/\sum|y|$ es un error relativo agregado: los clientes grandes dominan y no usa un baseline; además es robusta a $y\approx0$.
\end{itemize}
\endgroup

\section{Protocolo y métricas}
Una sola función \texttt{backtest(modelo)}: 92 orígenes diarios a las 00:00 (2014-10-01 a 2014-12-31); cada origen usa
datos hasta las 23:00 anteriores. Métricas en kWh: MAE, RMSE, WAPE ($\sum|e|/\sum|y|$), MAPE y MASE (escala: MAE del
naive estacional de 24 h en entrenamiento, por cliente). El MAPE se dispara con consumo cercano a cero (división entre $\approx0$).

\textbf{Referencia (seasonal naive 24 h, prueba):} MAE 49.4, RMSE 391.8, WAPE 8.47\%, MAPE 14.97\%, MASE 1.046;
MAE en la semana 22--28 dic: 100.4 (el doble, por el cambio de régimen navideño).

\section{Modelos, iteraciones y resultados}
\todo{tabla de iteraciones (L, meses de historia, hiperparámetros, pérdidas, MAE/MASE val, parámetros, tiempos, memoria GPU)
y curvas de pérdida de $\ge3$ iteraciones}

\todo{comparación de estrategias multi-paso del LSTM (autorregresiva, multi-salida, encoder--decoder) y MAE por $h$}

\todo{tabla comparativa ordenada por MASE: familia, L, estrategia, covariables, parámetros, tiempos, hardware,
MAE/RMSE/WAPE/MASE (H=24, H=168), MASE 22--28 dic, pinball y cobertura P10--P90}

\todo{3 gráficas: MASE vs tiempo de entrenamiento (Pareto), MAE por paso del horizonte, MASE vs meses de historia}

\todo{local vs global (10 clientes)}

\section{Discusión y conclusiones}
\todo{responder las 8 preguntas de la sección 8 del enunciado y desviación estándar del MASE con 3 semillas}

\bigskip
\noindent\textbf{Repositorio:} \todo{enlace al repositorio Git}
\end{document}
